% 本文件由 LaTeX 排版 Agent 根据有效 translation.json 自动生成。
% author=Apocrypha; work_id=0853; title=宗徒阿戴的训诲

\chapter{宗徒\Person{阿戴}的训诲}

% structure: p0001 source=h1 target=omitted reason=page-title-already-rendered-by-parent

\Person{阿戴}对他说：因为你这样相信了，我就因你所相信的那位的名，按手在你身上。就在他按手在他身上的那一刻，他长久所患的疾病之灾便得了痊愈。\Person{阿布加}感到惊奇而诧异，因为正如他所听到的关于\Person{耶稣}如何行事和医治，同样地，\Person{阿戴}也不用任何药物，就因\Person{耶稣}的名施行医治。\Person{阿布杜}之子\Person{阿布杜}脚患痛风，他也把脚伸到他面前，他按手在脚上，治好了他，他不再有痛风。在全城里，他也行了伟大的医治，显出奇妙的大能作为。\par

\Person{阿布加}对他说：如今人人都知道，你是藉\Person{耶稣基督}的权能行这些奇迹，看哪，我们为你的作为感到惊奇，因此我恳求你向我们讲述关于基督来临的事迹，它是如何发生的，以及祂光荣的大能，还有我们所听见祂所行的奇迹，这些你与你的同门都亲眼见过。\par

\Person{阿戴}说：我决不会闭口不言此事；因为我被差遣到这里，正是为此目的，要向每一个愿意相信的人说话和教导，就像你一样。明天请把全城的人聚集到我这里，我要借着向你们宣讲的话语，将生命之言播种在其中——关于基督的来临，它是如何发生的；关于那差遣祂的那一位，为何以及如何差遣祂；关于祂的大能和祂奇妙的作为；关于祂来临的光荣奥秘，就是祂在世上所谈论的；关于祂宣讲的准确真理；以及祂如何为了何事而自贬，并藉所取的人性贬抑祂崇高的天主性，被钉十字架，下到死者之处，冲破了那从未被冲破的围栏（\BibleRef{弗 2:14}），并藉自己被杀害而赐生命给死者，独自下去，却带着许多人上升到祂光荣的圣父那里，祂从永远就与圣父同在同一个崇高的天主性中。\par

\Person{阿布加}吩咐人把金银给\Person{阿戴}。\Person{阿戴}对他说：我们怎能收取不属于我们的东西呢？因为看哪，我们已经舍弃了属于我们的，正如我们的主所吩咐的；因为我们没有钱囊，没有口袋，将十字架扛在肩上，奉命在普天下宣讲祂的福音，祂为了我们被钉十字架，为救赎众人，整个受造界都感受到并为此受苦。\par

他在\Person{阿布加}王面前，在他的王侯贵族面前，在\Person{阿布加}的母亲\Person{奥古斯丁}面前，在\Person{阿布加}的妻子、\Person{默赫尔达特}之女\Person{沙尔玛特}面前，讲述了我们主的记号、祂的奇事、祂所行的光荣大能作为、祂的神性功业，以及祂升到圣父那里去；并讲述他们如何在祂被接升天时领受了权能和权柄——藉着同一权能，他治好了\Person{阿布加}和\Person{阿布杜}之子\Person{阿布杜}，他王国的第二人；以及祂如何告知他们，祂要在时代的终结和万物的终末显现自己；还有那将要来临的、为众人而有的苏醒与复活，以及绵羊与山羊之间、信徒与不信者之间将要进行的分离。\par

他对他们说：因为生命之门狭窄，真理之路窄小，所以相信真理的人很少，而由于不信，撒殚得着满足。因此说谎的人很多，他们误导那些能看见的人。因为，若不是有美好的结局等待着相信的人，我们的主就不会从天上降下来，来受生，并忍受死亡的痛苦。然而祂确实来了，并差遣了我们……关于我们所宣讲的信德：天主为众人被钉十字架。\par

如果有人不愿意同意我们这些话，就让他们到我们跟前来，向我们说明他们的心意，这样我们就可以像对待疾病一样，将医治的药应用于他们的思想，以治愈他们的病患。因为，虽然你们没有在基督受苦的时候在场，但藉着那变暗的太阳——你们看见了它——你们可以学习并明白那时所发生的大震动，当祂被钉十字架时，祂的福音藉着祂的同伴——我的同门——所行的记号传遍了全地；是的，那些希伯来人，只懂得他们生来就说的希伯来语，看哪，如今正说着各种语言，好让远处的人也能听见而相信，正如近处的人一样。因为正是祂混乱了在这地区我们之前那些傲慢者的口舌，也正是祂今天藉着我们这些来自巴勒斯坦\Place{加里肋亚}的卑微可鄙之人，教导真理和确切的信德。因为我——你们所看见的我——来自\Place{帕内阿斯}，就是约但河发源之地，我和同伴们被拣选作宣讲者。\par

因为，正如我的主所吩咐我的，看哪，我宣讲并传扬福音，看哪，我把祂的银钱抛在你们面前的桌子上，把祂话语的种子播种在众人的耳中；那些愿意接受的人，宣认基督的美好报偿就是他们的；但那些不信服的人，我就把脚上的尘土跺下，反对他们，正如祂所吩咐我的。\par

所以，我亲爱的，你们要悔改邪恶的道路和可憎的行为，以良善诚实的心志转向祂，正如祂以丰盛慈悲的恩宠转向你们；不要像从前逝去的世代，他们因硬着心对抗敬畏天主，公开受罚，好使他们自己受管教，也让后来的人战栗害怕。因为我们主来到世上的目的，确是要教导我们、指示我们：在创造的终结时，将有一切人的复活，那时他们的行为举止将显现在他们身上，他们的身体将成为正义书写的卷册；那里没有不识字的人，因为每个人都必读自己书上所写的。\par

你们有眼，却不觉察，自己也成了那看不见、听不见的人；你们无效的声音徒然向聋耳呼喊。虽然\Emph{他们}因听不见而无可责备，因为他们生来聋哑，但应得的责备却落在你们身上，因为你们不愿察觉——甚至不愿察觉你们所看见的。因为弥漫你们心智的谬误暗云，使你们无法获得天上的光明，即知识之领悟。\par

因此，逃离那些被造之物，如我告诉你们的，它们只是名义上被称为神，本性却非神；接近这位\Emph{存有}，祂按其本性是从永远到永远的天主，不是像你们的偶像那样被造的，也不是受造物和工艺品，如你们所夸耀的那些像。因为，虽然后者\Emph{存有}穿上身体，但祂\Emph{仍然}是与圣父同等的天主。因为受造之工——当祂被杀时颤抖，因祂受死而惊惶——这些见证祂就是创造主天主。因为大地颤抖不是因一个人，而是因那在水上建立大地者；太阳在天空变暗不是因一个人，而是因那造大光者；义人复活也不是因一个人，而是因那从起初就掌管死亡权势者；犹太圣所幔子从上到下裂开也不是因一个人，而是因那对他们说：\BibleQuote{你们的房屋将给你们留下一片荒凉}。因为，看！除非钉死祂的人知道祂是天主子，他们就不会宣告自己城的荒凉，也不会把祸患引到自己身上。因为，即使他们想轻看这宣认，当时发生的可怕震动也不容他们如此。因为，看！甚至钉死者的子女，如今也成了宣道者和福音传播者，与我的同宗徒们一起，在全巴勒斯坦地、在撒玛黎雅人中、在培肋舍特全地。异教的偶像也被藐视，基督的十字架受尊崇，\Emph{万国}万民都宣认降生成人的天主。\par

所以，如果在我们主耶稣尚在世上时，你们就相信祂是天主子，在听到祂宣讲的话之前，就已宣认祂是天主；现在祂已升到圣父那里，你们看见了因祂之名所行的神迹奇事，亲耳听到了祂福音的话，那么你们中间不可有人心里怀疑——好叫祂赐给你们的祝福应许得以在你们身上实现：那些没有看见我而相信我的，是有福的；而且，因为你们这样相信了我，你们所住的城将要蒙福，仇敌永远不能胜过它。所以，不要背离祂的信德。因为，看！你们已经听见看见了那些为祂的信德作证的事——\Emph{表明}祂是可敬爱的圣子，是荣耀的天主，是得胜的君王，是大能的权能；藉着对祂的信德，人能获得真正心智的眼睛，并理解：凡敬拜受造物的，正义的义怒必将临到他。\par

因为\Emph{在}我们向你们所说的每件事上，都是按照我们从主所领受的恩赐，我们就\Emph{这样}说、教导并宣讲\Emph{它}，为要使你们获得救恩，不致因异教的错误而毁灭自己的灵魂。因为天上的光明已照耀在受造界上，祂就是那拣选了古时的先祖、义人和先知，并藉圣神的启示与他们说话的那一位。因为祂就是钉死祂的犹太人的天主；而那些误入歧途的异教徒，即使不知道，也是在向祂献上敬拜。因为在天上地下没有别的神。看！从受造界的四方，宣认上达于祂。看！因此，你们的耳朵听到了你们未曾听过的；看！此外，你们的眼睛看到了你们未曾见过的。\par

所以，不要作你们所见所闻的反对者。当除去你们父辈的悖逆之心，从罪恶的轭下解放自己，这罪恶藉着在雕刻的像前奠酒和祭献\Emph{奉献}统治你们；要关切你们濒危的救恩，以及你们所倚靠的无效支持；要获得一颗新心，敬拜造物主而非受造之物——\Emph{这心里}描绘着真与实的肖像，就是圣父、圣子、圣神的肖像；相信并受洗于那三重荣耀的名号中。因为这就是我们的教导和宣讲。因为对基督真理的信仰并不包含许多事。你们中那些愿意服从基督的人，知道我在你们面前多次重复我的话，好叫你们学习并理解你们所听见的。\par

我们自己也必因此喜乐，如同农夫因田地蒙福而喜乐；天主也必因你们悔改归向祂而受光荣。你们由此得救，我们这些给你们这劝告的人，也不会丧失这工作的蒙福赏报。因为我确信你们是按主基督的旨意蒙福的土地，因此，代替我们奉命要向不接受我们言语的城镇抖下的脚上的尘土，我今天在你们耳门抖下了我口中的言论，其中描绘了基督\Emph{已经}来临和\Emph{将要}来临的事；以及复活、众人的苏醒、信者与不信者之间的分离；不认识天主的人所保留的严厉惩罚，以及那些相信基督、敬拜祂和祂崇高圣父、并宣认祂和祂神圣圣神者所将获得的未来福乐之蒙福应许。\par

如今，我应当结束当前的讲论；让那些接受了基督圣言的人与我们同留，也请那些愿意与我们一同祈祷的人留下；然后让他们各自回家。\par

宗徒阿戴乌斯见城中大批居民与他同留，心中欢喜；那时\Emph{仅仅}少数人没有留下，而就连那少数人在不多几日后也接受了他的言语，相信了在基督宣讲中所传的福音。\par

当宗徒阿戴乌斯在厄德撒全城面前讲了这些话，阿布加尔王看见全城男女老少都因他的教导而喜乐，\Emph{并听见}他们对他说：\Quote{打发你到我们这里来的基督是真实信实的}——他自己也为此大大喜乐，赞美天主；因为正如他从他的书记哈南那里听到关于基督的事，他也看见了宗徒阿戴乌斯因基督之名所行的奇妙大能作为。\par

阿布加尔王也对他说：正如我写信给基督所表明的，也正如基督给我回信所说的，\Emph{照样}我今天也从你本人领受了；所以我一生一世都要相信，并且持守同一件事而夸耀，因为我也知道：这些征兆和奇事除了因你们真诚真实所宣讲的基督的大能之外，没有别的能力能行。从今以后，我要敬拜祂——我和我的儿子马努、奥古斯丁，以及王后沙尔玛特。现在，随你愿意在哪里建造一座教堂，为那些相信并将相信你言语的人提供聚会的场所；并照你的主所赐给你的命令，按时放胆履行职务；对于将与你同作这福音教师的人，我准备给予丰厚的捐赠，使他们除了圣职之外别无它事；凡建造所需的一切费用，我必无限制地提供给你；你的言语在这城中将具有权威和主权；此外，\Emph{无需}他人介入，你可作为掌权者进入我王宫的殿宇。\par

阿布加尔下到他王宫后，他与他的贵族——阿布杜之子阿布杜、加尔迈、舍玛什格兰、阿布拜、梅赫尔达特，以及他们的同伴——都因他们亲眼所见、亲耳所闻的一切而喜乐；在他们内心的喜乐中，他们也开始赞美天主，因为天主使他们的心转向祂，弃绝了他们以往所生活的异教，宣认了基督的福音。阿戴乌斯建造了一座教堂后，他们和城中百姓在堂内献上誓愿和祭品；他们终身在那里不断献上赞美。\par

阿维达和巴尔卡尔巴，他们是首领和官长，佩戴王室的头带，走近阿戴乌斯，向他询问关于基督的事，\Emph{请求}他告诉他们：祂既是天主，又如何向他们显现为人？他们说：你们怎能看见祂？于是他就此事——凡他们亲眼所见、亲耳所闻的一切——完全满足了他们。此外，先知们论及祂的一切话，他都在他们面前重述，他们欢欣而满怀信心地接受了他的言语，没有一人反对他；因为他所行的荣耀事迹不容任何人反对他。\par

此外，沙维达和埃贝德内布，这城祭司的首领，连同他们的同伴皮罗兹和迪尔苏，当他们看见他所行的征兆时，就跑去推倒了他们素常向他们的神乃布和贝耳献祭的祭坛，唯有城中间的巨坛除外；他们呼喊说：这确实是那位杰出荣耀的主的弟子，我们曾听说祂在巴勒斯坦地所行的一切。凡相信基督的人，阿戴乌斯都接纳，以父、及子、及圣神之名给他们授洗。那些拜石头偶像的人坐在他脚前，从困扰他们的异教疯狂中得到康复。犹太人，经营细麻布的商人，熟悉法律和先知——他们也信服了，成为门徒，并宣认基督是永生天主之子。\par

但阿布加尔王和宗徒阿戴乌斯都没有用强力强迫任何人相信基督；因为不用人的强力，征兆的强力就迫使许多人相信了祂。他们满心爱慕地接受了他的教导——整个美索不达米亚地区和周围所有地区都是如此。\par

此外，制造王丝绸和头巾的阿盖乌斯，以及帕鲁特、巴尔舍拉马、巴尔撒米亚，连同其他同伴，都依附于宗徒阿代乌斯；他接纳了他们，并让他们与他一同从事职事，他们的工作就是阅读旧约和新约、先知书以及\BibleBook{}{宗徒大事录}，并每日默想它们；他严格命令他们，要使自己的身体纯洁、本人圣洁，如同在\Person{天主}的祭台前侍立的人所当行的。\Quote{你们要远离，}他说，\Quote{远离虚假誓言和邪恶杀人，以及不诚实的见证——这与奸淫相联；远离毫无怜悯的巫术，以及占卜、算命和卜卦者；远离命运和命理，这些是受迷惑的加色丁人所夸耀的；远离星辰和十二宫，愚昧人信赖这些。你们要远离不义的偏私、贿赂和馈赠，因为这些使无辜者被判有罪。与你们蒙召参与的这职事一起，你们不可有其他的工作；因为你们的职事工作就是\Person{主}，终身如此。你们要勤勉地授予圣洗圣事。不可贪爱现世的利益。要以正义和真理听讼。不可成为盲人的绊脚石，免得通过你们，那开了盲人眼睛者的名号被亵渎——正如我们所见的。愿所有看见你们的人，都察觉到你们自身与你们所宣讲和教导的一切和谐一致。}\par

他们在阿代乌斯奉王阿布加尔的命令和指示所建造的教会中与他一同服务，由国王及其贵族提供供给，一部分用于\Person{天主}的殿宇，一部分用于赒济穷人。此外，每天有许多人聚集并前来参加侍奉的祈祷，以及阅读旧约和\WorkTitle{四福音合参}的新约。他们也相信死者复活，并在复活希望的期盼中埋葬去世的人。他们按季节遵守教会的庆节，每日在教会的守夜礼中勤奋不懈。他们进行施舍访问，探望病人和健康的人，遵照阿代乌斯给他们的指示。在城市周边也建造了教会，许多人从他手中接受了司铎圣秩。甚至东方的民众，以商人装扮，进入罗马人地区，以便观看阿代乌斯所行的奇迹。那些成为门徒的人从他手中接受了司铎圣秩，并在他们本国的亚述人中教导自己的同胞，在那里秘密地建造祈祷之所，因为存在来自崇拜火、崇敬水之人的危险。\par

此外，亚述王纳尔塞斯听到宗徒阿代乌斯所行的同样事迹，便给王阿布加尔送去使命：要么把你面前行这些奇迹的人交给我，让我见见他并听他的言词；要么把你在你城内所见他所行的一切报告给我。于是阿布加尔写信给纳尔塞斯，向他从头至尾叙述了阿代乌斯所有事迹的整个故事；他什么也没有遗漏，全都写信告诉了他。纳尔塞斯听到写给他的那些事，极为惊讶和诧异。\par

此外，王阿布加尔因为不能前往罗马人地区，到巴勒斯坦去杀那些钉死基督的犹太人，便写了一封信，送交给提庇留凯撒，信中这样写道：——\par

王阿布加尔致我主提庇留凯撒：虽然我知道没有什么事能向陛下隐瞒，但我写信通知您可畏而强大的王权：在您统治之下、居住在巴勒斯坦地区的犹太人聚集起来，钉死了基督，祂没有任何当死的过犯，而在他们面前行了迹象和奇事，向他们显示了强大的大能作为，甚至为他们使死人复活；在他们钉祂时，太阳变黑，大地也震动，所有受造物都颤抖震动，仿佛出于自身，因这行为，整个受造界及其居民都退缩。如今陛下知道，对于做了这些事的犹太民族，您应下令如何处置。\par

提庇留凯撒写信给王阿布加尔，他这样写道：——\par

我收到了你对我忠诚的信件，并在我的面前诵读。关于犹太人在十字架事上胆敢做的，总督比拉多也写信通知了我的总督奥比努斯，关于你写信给我的同样事情。但是，因为此时正与反叛我的西班牙人民进行战争，故此我还未能为这事伸张正义；但我已准备好，当我有闲暇时，要依法对不依法行事的犹太人发出命令。并且为此，关于我任命为那里总督的比拉多——我已另派一人接替他，并因他背离法律、服从犹太人的意愿、为取悦犹太人而钉死基督，将他羞辱地免职；而依照我所听到的关于祂的，基督本应受他们尊敬和崇拜，而非遭受十字架的死亡；更因为他们亲眼看见了祂所行的一切。至于你，按照你对我的忠诚，以及你和你祖先所立下的忠心盟约，你很好地这样写信给我。\par

\Person{阿布加尔}王接待了由\Person{提庇留}凯撒派来的\Person{阿里斯提德}，并\Emph{回赠}了适合派他来者身份的尊贵礼物。\Person{阿里斯提德}从\Place{埃德萨}前往\Place{提昆塔}，那里有皇帝的副手\Person{克劳狄}；又从那里前往\Place{阿提卡}，\Person{提庇留}凯撒所在之处；\Person{盖乌斯}则负责守卫凯撒周围的区域。\Person{阿里斯提德}亲自在\Person{提庇留}面前讲述了\Person{阿达约}在\Person{阿布加尔}王面前所行的奇事。当战事稍闲，\Person{提庇留}便派人处死了\Place{巴勒斯坦}犹太人中一些首领。\Person{阿布加尔}王听说此事，极为欢喜，因为犹太人受了应得的惩罚。\par

\Person{阿达约}宗徒在\Place{埃德萨}建造了教堂，并为其配备了所需的一切，使城中大量居民成为门徒。此后数年，他又在\Emph{远近}的村庄建造教堂，将其完成并装饰，在各地设立执事和长老，教导读经的人，讲授内外礼仪与职事。\par

这一切之后，他染上疾病，即将离世。他在全体教会会众面前召叫\Person{阿加约}，命他近前，立他为继任的导师与治理者。原本是执事的\Person{帕鲁特}被立为长老；原本是书记的\Person{阿布舍拉玛}被立为执事。贵族与首领们聚集站立在他身旁——\Person{巴尔卡尔巴}（\Person{扎提}之子）、\Person{玛利亚哈布}（\Person{巴尔舍玛什}之子）、\Person{塞纳克}（\Person{阿维达}之子）、\Person{皮罗兹}（\Person{帕特里克}之子）以及其余同伴——\Person{阿达约}宗徒对他们说：\par

你们这些听我说话的人都知道且作见证：我素日在你们中间，凡我所宣讲、教导、你们从我所听见的，我都\Emph{照样}行在你们中间，你们也亲眼在行为上\Emph{看见了}。因为我们的主这样命令我们：我们在人前用言语宣讲什么，就当在众人面前用行为践行。你们也要按照\Place{耶路撒冷}门徒所订立、我的同僚宗徒们所遵行的典章和法律而行，不可偏离，不可删减。我亦在你们中间遵行这些，未尝偏左偏右，免得我疏远那应许的救恩——这救恩是为那些遵行的人所保留的。\par

因此，你们要留意所领受的职事，存着敬畏战兢的心恒守其中，每日履行。不可因疏忽习惯而履行，却要以信德的谨慎；赞美基督的话语不可从你们口中断绝，也不可对\Emph{定时}祈祷感到厌倦。要留意你们所持守的真理、所领受的真道教训，以及我托付给你们的救恩的产业，因为你们将在基督的审判台前为此交账——当祂与牧者及监督算账，并从商人那里连本带利收回祂的钱财时。因为祂是君王之子，要去接受王国并返回；祂必再来，使众人复活\Emph{得生命}，然后坐在祂正义的宝座上，审判死者与活人，正如祂对我们所说的。\par

不可让骄傲遮蔽你们内心的隐秘眼睛，以免你们在原本无绊脚石的路上绊倒许多，反而使那路上的迷途可憎。你们要寻找迷失的，引领迷途的，为找到的而喜乐；包扎受伤的，看顾肥壮的，因为基督的羊群将在你们手中被追讨。不要贪图那消逝的荣耀：那期待从羊群得荣耀的牧人，他的羊群与他相处何等可悲。你们要格外关心年幼的羔羊，他们的天使常瞻仰那不可见的圣父的容颜。你们在犹太人——钉十字架者，以及受迷惑的外邦人中，不可成为盲人的绊脚石，却要成为崎岖之地的开路者与修路者。因为你们需与这两方争战，以显明你们所持守的信仰真理；即使你们沉默，你们谦逊端正的外表也将为你们向那恨真理、爱谎言的人争战。\par

不可在富人面前责打穷人，因为贫穷本身已是他们沉重的鞭笞。\par

\Quote{不可被撒殚的可憎诡计所迷惑，免得你们被剥去所穿戴的信仰。……\InnerQuote{我们与那钉十字架的犹太人毫无往来。我们从你领受的这产业，我们必不放弃，必要持守它离世；在我们主的日子，在祂正义的审判座前，祂必将这产业归还我们，正如你所告诉我们的。}}\par

这些话说完后，\Person{阿布加尔}王与他的首领和贵族起身回宫，众人皆因他垂死而忧伤。他派人送去高贵华美的寿衣，好让\Person{阿达约}穿着下葬。\Person{阿达约}见后，派人回复\Emph{说}：\Quote{我活着时未曾收取你的任何东西，现在临终也绝不收取你的任何东西，也不愿违背基督对我们所说的话：\InnerQuote{不可从任何人收取什么，也不可在世上拥有什么。}}\par

在这些事由宗徒阿代宣讲之后又过了三日，他听到了并接受了关于他们宣讲中所陈明的教导的见证，与那些同他一起从事职务的人在一起，在全体显贵面前，他离开了这个世界。那一日是星期四，依雅尔月十四日，\Emph{大致相当于五月}。全城都为他极其悲伤，痛心疾首。不仅基督徒为他忧愁，连同在这城中的犹太人和外邦人也如此。但阿布加尔王比任何人更为他忧愁，他和他的王国中的公侯们。在他灵魂的悲伤中，他在那一日鄙弃并搁置了他王权的辉煌，以夹杂着哀叹的眼泪同众人一起哀悼他。全城所有见到他的人都惊讶地\Emph{看到}他为他所受的巨大痛苦。他以盛大无与伦比的排场\Emph{抬着}他，像安葬一位去世的公侯一样埋葬了他；他将他安放在一座精美雕刻装饰的宏伟墓穴中——那是安放阿布加尔王祖先阿留家族之人的墓穴：他在忧伤、悲哀和极大的痛苦中将他安放在那里。教会的全体民众不时前往那里热切祈祷；他们年复一年地纪念他的离世，遵照宗徒阿代所交代给他们的命令和指示，以及阿盖乌斯的话——他本人成为了导师与管理者，与他之后的继承者，通过他所接受的、在众人面前由他授予的司铎圣秩而接替他的职位。\par

他也以从他那里接受的同一圣秩，在美索不达米亚整个地区任命了司铎和导师。因为他们也如同宗徒阿代一样，持守他的话，聆听并接受\Emph{它}，作为可敬基督的宗徒的良好而忠实的继承人。但他不接受任何人的金银，王公们的礼物也未曾接近他：因为，\Emph{接受}金银，他\Emph{自己}反以信徒的灵魂丰富了基督的教会。\par

此外，\Emph{至于}男人和女人的整个状态，他们是贞洁而谨慎的，神圣而纯洁的：因为他们如同隐修士一般贞洁地生活，毫无玷污——在职务上谨慎守望，同情穷人，探访病人：因为他们的脚步充满了所见之人的赞美，他们的行为身着陌生人嘉许的外衣——甚至乃步和贝耳家中的司祭也时时与他们分享尊荣，由于他们庄重的神态，诚实的言语，因高贵天性而来的坦率言辞，既不因贪婪而屈从，也不因\Emph{惧怕}责备而受束缚。因为没有人见到他们而不跑上前去恭敬地请安，因为仅见到他们就将平安倾注在观看者身上：因为他们的温柔言语就像一张网铺在倔强之人身上，他们进入了真理和真实的圈中。因为没有人见了他们而为他们羞愧，因为他们没有做任何不合乎正直与得体的事。因此，他们的举止无所畏惧，向众人宣讲他们的教导。因为凡对别人所说并吩咐的，他们自己亲身在实践中体现出来；而听者看到他们的行为与言语一致，无需太多说服就成为他们的门徒，并承认基督为君王，赞美天主使他们归向祂。\par

阿布加尔王死后数年，他有一个倔强的儿子兴起，不喜和平；他\Emph{传话}给阿盖乌斯，当时他正坐在教会中：\Quote{为我做一条金头带，如同你从前为我父辈们所做的那样。}阿盖乌斯回复他：\Quote{我不愿放弃基督交托给我的职务，而去制作一条邪恶的头带。}当他看到他不顺从时，便派人趁他坐在教会中讲解时打断了他的腿。临终时，他嘱咐帕鲁特和阿布舍拉玛：\Quote{在这座因真理之故，看哪！我正死去的房屋里，将我安放埋葬。}他们便照他所嘱咐的，将他安放在教会中间的门口，在男人和女人之间。全教会和全城都极其悲痛哀悼——除了教会中已有的痛苦和哀悼之外，如同宗徒阿代本人去世时的哀悼一样。\par

由于他在断腿时突然迅速去世，他未能将手按在帕鲁特身上。帕鲁特前往安提约基雅，从安提约基雅主教塞拉比翁接受了司铎圣秩；此塞拉比翁本人也从罗马城的主教则斐琳接受了圣秩，这是承继自西满·刻法之司铎圣秩的传承；西满·刻法从我主领受了圣秩，并在罗马作主教二十五年，当时在位的是执政十三年的凯撒。\par

依照阿布加尔王国内及所有王国中存在的惯例，凡是王所命令和凡在他面前所说的话，都被记载并存放于档案中；同样，塞纳克之子、厄贝德沙代之孙、国王的书记官拉布纳，从始至终也将这些有关宗徒阿代的事记载下来，同时王室的文书哈南——一名国王的文书——也以见证之笔附署，并将\Emph{文件}存放在国王的档案中，那里存放着法令和法律，以及买卖\Emph{契约}被小心保存，没有任何疏忽。\par

\Emph{此处}结束了宗徒阿代的教导，即他在厄德萨——信实的阿布加尔王的信实之城——所宣讲的教导。\par

\Emph{译者注}——下列叙利亚月份名称表，这些名称在帝国及塞琉古时代使用，其中多个已在这些文献中出现，取自Volck于1859年编辑的阿拉伯文和拉丁文本\WorkTitle{Caswinii Calendarium Syriacum}。此处另附希伯来文名称以供比较。但须注意，\Quote{叙利亚历法中使用的年份，至少在降生成人后，是儒略年，由罗马月份组成。}（参见\WorkTitle{L'Art de vérifier les dates: }巴黎，1818年，卷I，第45页。）因此，其与希伯来月份的对应并不如名称所示那样精确，因为希伯来月份始于新月，并在第十二个月阿达尔之后增加一个闰月——韦亚达尔。\par

\Emph{月份 / 叙利亚名 / 希伯来名}\par

十月 / Tishri prior / Tishri, or Ethanim\par

十一月 / Tishri posterior / Bull, or Marcheshvan\par

十二月 / Canun prior / Chisleu\par

一月 / Canun posterior / Tebeth\par

二月 / Shubat / Shebat\par

三月 / Adar / Adar\par

四月 / Nisan / Nisan\par

五月 / Ajar / Zif, or Iyar\par

六月 / Chaziran / Sivan\par

七月 / Tamuz / Tammuz\par

八月 / Ab / Ab\par

九月 / Elul / Elul\par

\SourceRecord{Translated by B.P. Pratten.；Ante-Nicene Fathers；Vol. 8.；Edited by Alexander Roberts, James Donaldson, and A. Cleveland Coxe.；Buffalo, NY: Christian Literature Publishing Co.,；1886.；<http://www.newadvent.org/fathers/0853.htm>.}
